% Part: incompleteness
% Chapter: representability-in-q
% Section: sigma1-completeness

\documentclass[../../../include/open-logic-section]{subfiles}

\begin{document}

\olfileid{inc}{inp}{s1c}
\olsection{\texorpdfstring{$\Sigma_1$}{Sigma-1} کامل‌بودن}

با وجود ناتمامیت~$\Th{Q}$ و گسترش‌های سازگار و اصل‌موضوعه‌پذیر آن،
دیدیم که~$\Th{Q}$ بسیاری از واقعیت‌های پایه دربارهٔ عددنماها را ثابت
می‌کند. در واقع، می‌توان این نتیجه را به‌طور چشمگیری گسترش داد. برای
درک گسترهٔ آنچه در~$\Th{Q}$ اثبات‌پذیر است، مفاهیم !!{formula}sی
$\Delta_0$، $\Sigma_1$ و~$\Pi_1$ را معرفی می‌کنیم. به‌طور تقریبی،
!!{formula} $\Sigma_1$ فرمولی به صورت~$\lexists[x][!B(x)]$ است که
در آن~$!B$ تنها با پیوندهای گزاره‌ای و سورهای کراندار ساخته شده است.
نشان خواهیم داد اگر~$!A$ یک !!{sentence} $\Sigma_1$ باشد که در
$\Struct{N}$ صادق است، آنگاه~$\Th{Q} \Proves !A$
(\olref{thm:sigma1-completeness}).

\begin{defn}
\ollabel{defn:bd-quant}
یک \emph{!!{formula} وجودی کراندار} فرمولی به صورت
$\lexists[x][(x < t \land !A(x))]$ است، که در آن~$t$ هر ترمی است؛
آن را به‌طور قراردادی به صورت~$\bexists{x < t}{!A(x)}$ می‌نویسیم.
%
یک \emph{!!{formula} کلی کراندار} فرمولی به صورت
$\lforall[x][(x < t \lif !A(x))]$ است، که در آن~$t$ هر ترمی است؛
آن را به‌طور قراردادی به صورت~$\bforall{x < t}{!A(x)}$ می‌نویسیم.
\end{defn}

\begin{defn}
\ollabel{defn:delta0-sigma1-pi1-frm}
!!{formula}~$!B$ را $\Delta_0$ می‌نامیم اگر از !!{formula}sی اتمی،
تنها با استفاده از پیوندهای گزاره‌ای و سورگذاری کراندار ساخته شده باشد.
%
!!{formula}~$!A$ را $\Sigma_1$ می‌نامیم اگر
$!A \ident \lexists[x][!B(x)]$، که در آن~$!B$، $\Delta_0$ است.
%
!!{formula}~$!A$ را $\Pi_1$ می‌نامیم اگر
$!A \ident \lforall[x][!B(x)]$، که در آن~$!B$، $\Delta_0$ است.
\end{defn}

\begin{lem}
\ollabel{lem:q-proves-clterm-id} فرض کنید~$t$ ترمی بسته باشد به‌گونه‌ای
که~$\Value{t}{N} = n$. آنگاه~$\Th{Q} \Proves \eq[t][\num n]$.
\end{lem}

\begin{proof}
این را با استقرا روی پیچیدگی~$t$ ثابت می‌کنیم. در حالت پایه،
$\Value{\Obj 0}{N} = 0$ و~$\Th{Q} \Proves \eq[\Obj 0][\num 0]$،
زیرا~$\num 0 \ident \Obj 0$.
%
برای حالت استقرایی، بگذارید~$t_1$ و~$t_2$ ترم‌هایی باشند به‌گونه‌ای که
$\Value{t_1}{N} = n_1$، $\Value{t_2}{N} = n_2$،
$\Th{Q} \Proves \eq[t_1][\num n_1]$ و
$\Th{Q} \Proves \eq[t_2][\num n_2]$.

آنگاه~$\Value{(t_1')}{N} = n_1 + 1$، و داریم
$\Th{Q} \Proves \eq[t_1'][{\num n_1}']$؛ این از اعمال قواعد
مرتبه‌اول این‌همانی به فرض استقرا و !!{formula}
$\eq[\num{n_1}'][\num{n_1}']$ به دست می‌آید. پس بنا بر تعریف عددنماها
$\Th{Q} \Proves \eq[t_1'][\num{n_1 + 1}]$.

برای جمع‌ها داریم
\[
      \Value{(t_1 + t_2)}{N}
    = \Value{t_1}{N} + \Value{t_2}{N}
    = n_1 + n_2.
\]
بنا بر فرض استقرا و قواعد این‌همانی،
$\Th{Q} \Proves \eq[t_1 + t_2][\num{n_1} + t_2]$، و سپس با کاربرد
دوبارهٔ قواعد این‌همانی
$\Th{Q} \Proves \eq[t_1 + t_2][\num{n_1} + \num{n_2}]$.
بنا بر~\olref[inc][req][bre]{lem:q-proves-add}،
$\Th{Q} \Proves \eq[\num{n_1} + \num{n_2}][\num{n_1 + n_2}]$؛
پس~$\Th{Q} \Proves \eq[t_1 + t_2][\num{n_1 + n_2}]$.

استدلالی مشابه برای~$\times$ نیز کار می‌کند، با استفاده از
\olref[inc][req][bre]{lem:q-proves-mult}.
%
چون بدین ترتیب همهٔ ترم‌های بستهٔ حساب را بررسی کرده‌ایم، داریم
$\Th{Q} \Proves \eq[t][\num n]$ برای همهٔ ترم‌های بستهٔ~$t$ که
$\Value{t}{N} = n$.
\end{proof}

\begin{prob}
بخش مربوط به~$\times$ از~\olref{lem:q-proves-clterm-id} را با جزئیات
ثابت کنید.
\end{prob}

\begin{lem}
\ollabel{lem:atomic-completeness}
فرض کنید~$t_1$ و~$t_2$ ترم‌هایی بسته باشند. آنگاه
\begin{enumerate}
\item اگر~$\Value{t_1}{N} = \Value{t_2}{N}$،
    آنگاه~$\Th{Q} \Proves \eq[t_1][t_2]$.
\item اگر~$\Value{t_1}{N} \neq \Value{t_2}{N}$،
    آنگاه~$\Th{Q} \Proves \eq/[t_1][t_2]$.
\item اگر~$\Value{t_1}{N} < \Value{t_2}{N}$،
    آنگاه~$\Th{Q} \Proves t_1 < t_2$.
\item اگر~$\Value{t_2}{N} \leq \Value{t_1}{N}$،
    آنگاه~$\Th{Q} \Proves \lnot(t_1 < t_2)$.
\end{enumerate}
\end{lem}

\begin{proof}
با داشتن ترم‌های~$t_1$ و~$t_2$، قرار می‌دهیم
$n = \Value{t_1}{N}$ و~$m = \Value{t_2}{N}$.

فرض کنید~$!A \ident t_1 = t_2$. بنا بر
\olref{lem:q-proves-clterm-id}،
$\Th{Q} \Proves \eq[t_1][\num n]$ و
$\Th{Q} \Proves \eq[t_2][\num m]$.
اگر~$n = m$، آنگاه~$\Th{Q} \Proves \eq[\num n][\num m]$ و بنابراین
به سبب تعدی این‌همانی~$\Th{Q} \Proves \eq[t_1][t_2]$.
اگر~$n \neq m$، آنگاه~$\Th{Q} \Proves \eq/[\num n][\num m]$، و
باز هم به سبب تعدی این‌همانی،
$\Th{Q} \Proves \eq/[t_1][t_2]$.

اکنون بگذارید~$!A \ident t_1 < t_2$. در هر دو حالت بر اصل موضوعهٔ
$!Q_8$ تکیه می‌کنیم، که می‌گوید
$x < y \liff \lexists[z][\eq[z' + x][y]]$ برای همهٔ~$x,y$.

فرض کنید~$\Sat{N}{t_1 < t_2}$. آنگاه~$k \in \Nat$ای وجود دارد
به‌گونه‌ای که~$n + k + 1 = m$. بنا بر
\olref{lem:q-proves-clterm-id}،
$\Th{Q} \Proves \eq[t_1][\num n]$ و
$\Th{Q} \Proves \eq[t_2][\num m]$، و بنا بر بخش نخست این لم،
$\Th{Q} \Proves \eq[{\num k}' + \num n][\num m]$.
از تعدی این‌همانی نتیجه می‌شود
$\Th{Q} \Proves \eq[{\num k}' + t_1][t_2]$؛ پس
$\Th{Q} \Proves \lexists[z][\eq[z' + t_1][t_2]]$.
از جهت راست‌به‌چپ~$!Q_8$، $\Th{Q} \Proves t_1 < t_2$.

در مقابل فرض کنید~$\Sat/{N}{t_1 < t_2}$، یعنی~$m \leq n$.
%
در~$\Th{Q}$ کار می‌کنیم و فرض می‌کنیم~$t_1 < t_2$. بنا بر جهت
چپ‌به‌راست~$!Q_8$، $z$ای وجود دارد به‌گونه‌ای که
$\eq[z' + t_1][t_2]$. چون~$\Th{Q} \Proves \eq[t_1][\num n]$ و
$\Th{Q} \Proves \eq[t_2][\num m]$، داریم
$\eq[z' + \num n][\num m]$.
%
با استقرایی بیرونی روی~$m$ و با استفاده از~$!Q_5$،
$\eq[z' + \num{n - m}][\Obj 0]$.
اگر~$m = n$، آنگاه~$\eq[z'][\Obj 0]$، که با~$!Q_2$ تناقض دارد.
اگر~$m < n$، آنگاه
$\eq[(z' + \num{n - m - 1})'][\Obj 0]$، که باز هم بنا بر~$!Q_5$
به دست می‌آید و با~$!Q_2$ تناقض دارد.
پس~$\Th{Q} \Proves \lnot(t_1 < t_2)$.
\end{proof}

\begin{lem}
\ollabel{lem:bounded-quant-equiv}
فرض کنید~$!A$ یک !!a{formula}، $t$ ترمی بسته و
$k=\Value{t}{N}$ باشد. آنگاه
\begin{enumerate}
\item $\Th{Q} \Proves \bforall{x<t}{!A(x)}$ اگر و تنها اگر
$\Th{Q} \Proves !A(\num 0) \land \dots \land !A(\num{k-1})$.
\item $\Th{Q} \Proves \bexists{x<t}{!A(x)}$ اگر و تنها اگر
$\Th{Q} \Proves !A(\num 0) \lor \dots \lor !A(\num{k-1})$.
\end{enumerate}
\end{lem}

\begin{proof}
    حالت سور کلی کراندار را ثابت می‌کنیم.
    اگر~$\Value{t}{N} = 0$، آنگاه سمت چپ هم‌ارزی در~$\Th{Q}$
    اثبات‌پذیر است، زیرا بنا بر
    \olref[inc][req][min]{lem:less-zero} هیچ~$x<\num 0$ای وجود ندارد.
    به همین ترتیب، می‌توانیم هم‌بندی تهی را صرفاً~$\ltrue$ بگیریم،
    که آن نیز در~$\Th{Q}$ اثبات‌پذیر است.
    %
    بنابراین فرض می‌کنیم~$\Value{t}{N} = k+1$ برای عدد طبیعی~$k$ای.
    بنا بر~\olref{lem:q-proves-clterm-id}
    می‌توانیم فرض کنیم با !!a{formula} به صورت
    $\bforall{x<\num{k+1}}{!A(x)}$ کار می‌کنیم.

    فرض کنید
    $\Th{Q} \Proves \bforall{x<\num{k+1}}{!A(x)}$، و بگذارید
    $n \leq k$. چون بنا بر~\olref{lem:atomic-completeness} داریم
    $\Th{Q} \Proves \num n < \num{k+1}$، از منطق نتیجه می‌شود
    $\Th{Q} \Proves !A(\num n)$. با به‌کاربردن این واقعیت~$k+1$
    بار، یکی برای هر~$n \leq k$، چنان‌که می‌خواستیم به
    $\Th{Q} \Proves !A(\num 0) \land \dots \land !A(\num k)$
    می‌رسیم.

    برای جهت دیگر، فرض کنید
    $\Th{Q} \Proves !A(\num 0) \land \dots \land !A(\num k)$.
    در~$\Th{Q}$ کار می‌کنیم و فرض می‌کنیم~$x < \num{k+1}$.
    بنا بر~\olref[inc][req][min]{lem:less-nsucc} داریم
    $x = \num 0 \lor \dots \lor x = \num k$؛ پس از منطق
    $!A(x)$ نتیجه می‌شود، و بنابراین ادعای کلی
    $\bforall{x<\num{k+1}}{!A(x)}$ به دست می‌آید.

    اثبات هم‌ارزی برای !!{formula}sی دارای سور وجودی کراندار مشابه است.
\end{proof}

\begin{prob}
حالت وجودی در~\olref{lem:bounded-quant-equiv} را با جزئیات ثابت کنید.
\end{prob}

\begin{lem}
\ollabel{lem:delta0-completeness}
اگر~$!A$ یک !!{sentence} $\Delta_0$ باشد که در~$\Struct{N}$ صادق است،
آنگاه~$\Th{Q} \Proves !A$.
\end{lem}

\begin{proof}
این را با استقرا روی پیچیدگی !!{formula} ثابت می‌کنیم.
%
حالت پایه را~\olref{lem:atomic-completeness} به دست می‌دهد، پس به
گام استقرا می‌رویم. برای سادگی، حالت نقیض را بر حسب ساختار
!!{formula}ی که نقیض بر آن اعمال شده است به زیرحالت‌هایی تقسیم می‌کنیم.

\begin{enumerate}
\item فرض کنید~$(!A \land !B)$ در~$\Struct{N}$ صادق باشد؛ پس~$!A$
و~$!B$ در~$\Struct{N}$ صادق‌اند. بنا بر فرض استقرا،
$\Th{Q} \Proves !A$ و~$\Th{Q} \Proves !B$؛ پس بنا بر منطق
$\Th{Q} \Proves (!A \land !B)$.
%
\item فرض کنید~$\lnot (!A \land !B)$ در~$\Struct{N}$ صادق باشد؛
پس یا~$\lnot !A$ یا~$\lnot !B$ در~$\Struct{N}$ صادق است. بدون کاستن
از کلیت، اولی را فرض کنید. بنا بر فرض استقرا
$\Th{Q} \Proves \lnot !A$، و بنابراین بنا بر منطق
$\Th{Q} \Proves \lnot (!A \land !B)$.
%
\item فرض کنید~$(!A \lor !B)$ در~$\Struct{N}$ صادق باشد؛ پس یا
$!A$ در~$\Struct{N}$ صادق است یا~$!B$ در~$\Struct{N}$. بدون کاستن از کلیت، فرض کنید
اولی برقرار است. بنا بر فرض استقرا~$\Th{Q} \Proves !A$، و بنابراین
بنا بر منطق~$\Th{Q} \Proves (!A \lor !B)$.
%
\item فرض کنید~$\lnot(!A \lor !B)$ در~$\Struct{N}$ صادق باشد؛ پس
$\lnot !A$ و~$\lnot !B$ در~$\Struct{N}$ صادق‌اند. آنگاه بنا بر فرض
استقرا~$\Th{Q} \Proves \lnot !A$ و
$\Th{Q} \Proves \lnot !B$. در نتیجه، بنا بر منطق
$\Th{Q} \Proves \lnot(!A \lor !B)$.
%
\item فرض کنید~$\bforall{x<t}{!A(x)}$ در~$\Struct{N}$ صادق باشد،
که در آن~$t$ ترمی بسته و~$k=\Value{t}{N}$ است. بنا بر فرض استقرا و
منطق، اگر~$!A(\num n)$ در~$\Struct{N}$ برای همهٔ
$n < \Value{t}{N}$ صادق باشد، آنگاه
$\Th{Q} \Proves !A(\num 0) \land \dots \land !A(\num{k-1})$.
از~\olref{lem:bounded-quant-equiv} نتیجه می‌شود
$\Th{Q} \Proves \bforall{x<t}{!A(x)}$.
%
\item حالت سور وجودی کراندار، که در آن !!a{sentence} به صورت
$\bexists{x < t}{!A(x)}$ داریم، مشابه حالت سور کلی کراندار است.
%
\item فرض کنید~$\lnot \bforall{x<t}{!A(x)}$ در~$\Struct{N}$ صادق
باشد، که در آن~$t$ ترمی بسته است. این !!{sentence} با !!{sentence}ٔ
$\bexists{x<t}{\lnot !A(x)}$ هم‌ارز است و هم‌ارزی در~$\Th{Q}$
اشتقاق‌پذیر است؛ پس می‌توانیم استدلال مربوط به سورهای وجودی کراندار را
به کار بریم.
%
\item به‌طور مشابه، فرض کنید~$\lnot \bexists{x<t}!A(x)$ در
$\Struct{N}$ صادق باشد، که در آن~$t$ ترمی بسته است. این !!{sentence} در
$\Th{Q}$ با~$\bforall{x<t}{\lnot!A(x)}$ هم‌ارز است؛ پس می‌توانیم
استدلال مربوط به سورهای کلی کراندار را به کار بریم.
%
\item سرانجام، فرض کنید~$\lnot !A$ در~$\Struct{N}$ صادق باشد. تنها
حالت‌های باقیمانده این‌اند که~$!A$ اتمی باشد یا
$\lnot !A \ident \lnot\lnot !B$ برای !!{sentence}
$\Delta_0$ای چون~$!B$. اگر~$!A$ اتمی باشد، بنا بر
\olref{lem:atomic-completeness}، $\Th{Q} \Proves \lnot !A$.
اگر~$\lnot !A \ident \lnot\lnot !B$، آنگاه بنا بر منطق در
$\Th{Q}$ به‌طور اثبات‌پذیر با~$!B$ هم‌ارز است؛ و آن در~$\Struct{N}$
صادق است، زیرا~$\lnot !A$ در~$\Struct{N}$ صادق است. بنابراین
بنا بر فرض استقرا~$\Th{Q} \Proves \lnot !A$.
\end{enumerate}
\end{proof}

\begin{prob}
حالت وجودی در~\olref{lem:delta0-completeness} را با جزئیات ثابت کنید.
\end{prob}

\begin{thm}
\ollabel{thm:sigma1-completeness}
اگر~$!A$ یک !!{sentence} $\Sigma_1$ باشد که در~$\Struct{N}$ صادق است،
آنگاه~$\Th{Q} \Proves !A$.
\end{thm}

\begin{proof}
اگر~$\lexists{x}!A(x)$ یک !!{sentence} $\Sigma_1$ باشد که در
$\Struct{N}$ صادق است، آنگاه عدد طبیعی~$n$ و انتساب متغیری چون~$s$
وجود دارند به‌گونه‌ای که~$s(x) = n$ و~$\Sat{N}{!A(x)}[s]$.
از واقعیت‌های استاندارد دربارهٔ رابطهٔ ارضا نتیجه می‌شود
$\Sat{N}{!A(\num n)}$. اما~$!A(\num n)$ یک !!{formula} $\Delta_0$
است؛ پس بنا بر~\olref{lem:delta0-completeness} داریم
$\Th{Q} \Proves !A(\num n)$، و بنابراین بنا بر منطق نیز
$\Th{Q} \Proves \lexists[x][!A(x)]$.
\end{proof}

\end{document}
